\section{Ethics and Privacy}

The facial study is optional and separable from participation in the core ARC-ADHD programme. Parent or guardian permission and child assent must explicitly cover camera processing. Declining the camera module should not remove access to the non-camera study when scientifically and operationally feasible.

Processing is performed on-device whenever possible. Raw video is not part of standard telemetry. Stored research data should preferentially contain derived temporal features, tracking confidence, and event-aligned summaries. Any exceptional retention of raw video for algorithm validation requires a separate justification, retention period, access policy, and consent language.

The GDPR defines biometric data in relation to technical processing that allows or confirms unique identification, and separately treats data concerning health as a special category \parencite{eu2016gdpr}. The study must therefore distinguish facial-behavior research features from identity-recognition biometrics while still recognizing that the broader research record can contain sensitive health information.

A data-protection impact assessment should explicitly map camera capture, transient processing, feature extraction, upload, storage, researcher access, withdrawal, and deletion. The adaptive system must never use facial behavior to infer identity, diagnose a disorder, or override an explicit child request or safety rule.
